\documentclass[tikz,border=3mm]{standalone}
\usepackage{listofitems}
\usetikzlibrary{positioning,calc,arrows.meta,shapes.misc}

% =========================
% Customization
% =========================
\newcommand{\CellLabel}{RNN Cell}
\newcommand{\LayerOneName}{1st Layer}
\newcommand{\LayerTwoName}{2nd Layer}
\newcommand{\LayerKName}{Kth Layer}
\newcommand{\PredictionLabel}{$\hat{y}$}

% Input labels, one entry per time step.
\readlist\TimeLabels{$u(t-k)$,$u(t-k-1)$,$u(t-k-2)$,$u(t-1)$}

% Spacing controls.
\def\StepSpacing{1.15cm}
\def\GapBeforeLastStep{2.0cm}
\def\LayerOneToTwoGap{1.5cm}
\def\LayerTwoToKGap{2.0cm}
\def\InputLabelGap{0.55cm}

% Colors.
\colorlet{LayerOneColor}{yellow!25}
\colorlet{LayerTwoColor}{blue!20}
\colorlet{LayerKColor}{green!22}

\tikzset{
  cell/.style={
    draw,
    rounded rectangle,
    minimum height=1.0cm,
    minimum width=1.7cm,
    font=\small
  },
  temporal/.style={->,>=Latex,line width=0.8pt},
  temporal skip/.style={temporal,dashed},
  vertical/.style={->,>=Latex,line width=0.8pt},
  vertical skip/.style={vertical,dashed}
}

\begin{document}
\begin{tikzpicture}
  \pgfmathtruncatemacro{\NumSteps}{\TimeLabelslen}
  \pgfmathtruncatemacro{\LastStep}{\NumSteps}
  \pgfmathtruncatemacro{\LastTransitionStart}{\NumSteps-1}

  % Base row (1st layer).
  \node[cell,fill=LayerOneColor] (L1-1) {\CellLabel};
  \foreach \Step in {2,...,\NumSteps} {
    \pgfmathtruncatemacro{\Prev}{\Step-1}
    \ifnum\Step=\NumSteps
      \node[cell,fill=LayerOneColor,right=\GapBeforeLastStep of L1-\Prev] (L1-\Step) {\CellLabel};
    \else
      \node[cell,fill=LayerOneColor,right=\StepSpacing of L1-\Prev] (L1-\Step) {\CellLabel};
    \fi
  }

  % 2nd and Kth layers.
  \foreach \Step in {1,...,\NumSteps} {
    \node[cell,fill=LayerTwoColor,above=\LayerOneToTwoGap of L1-\Step] (L2-\Step) {\CellLabel};
    \node[cell,fill=LayerKColor,above=\LayerTwoToKGap of L2-\Step] (LK-\Step) {\CellLabel};
  }

  % Horizontal temporal connections.
  \foreach \Step in {1,...,\LastTransitionStart} {
    \pgfmathtruncatemacro{\Next}{\Step+1}
    \ifnum\Next=\NumSteps
      \draw[temporal skip] (L1-\Step) -- (L1-\Next);
      \draw[temporal skip] (L2-\Step) -- (L2-\Next);
      \draw[temporal skip] (LK-\Step) -- (LK-\Next);
    \else
      \draw[temporal] (L1-\Step) -- (L1-\Next);
      \draw[temporal] (L2-\Step) -- (L2-\Next);
      \draw[temporal] (LK-\Step) -- (LK-\Next);
    \fi
  }

  % Vertical layer connections.
  \foreach \Step in {1,...,\NumSteps} {
    \draw[vertical] (L1-\Step) -- (L2-\Step);
    \draw[vertical skip] (L2-\Step) -- (LK-\Step);
  }

  % Layer labels.
  \node[left=0.6cm of L1-1,align=right] {\LayerOneName};
  \node[left=0.6cm of L2-1,align=right] {\LayerTwoName};
  \node[left=0.6cm of LK-1,align=right] {\LayerKName};

  % Input labels.
  \foreach \Step in {1,...,\NumSteps} {
    \node[below=\InputLabelGap of L1-\Step] {\TimeLabels[\Step]};
  }

  % Prediction arrow.
  \draw[vertical] (LK-\LastStep.north) -- ++(0,1.0) node[above] {\PredictionLabel};
\end{tikzpicture}
\end{document}
